%%%PAGE 112 rot=0 legibility=good summary=光学：衍射与干涉条件、双缝干涉、液体膜气体膜固体膜的薄膜干涉、单缝圆盘衍射、偏振
% 页面最上缘右上角露出的“Date. 混合气体 ρ=1.429 g·L^-1”“…=6.72 L”属相邻纸页（化学），按规范忽略
%%%BLOCK id=112-01 ch=15 sec=折射·全反射 kind=note alt=none cont=prev y=0.04-0.09
% DUP?: 页眉区（本页第一条横线之上）的字迹，其中“t d + En 同大小”与 p111 页底最后一行“(td+En) 同大小”相同，可能是相邻页字迹露出；若与 p111 重复可删除
\unclear{似保?}\quad \unclear{$t\,d\,+E_n$} 同大小
%%%END
%%%BLOCK id=112-02 ch=15 sec=光的衍射 kind=knowledge alt=none cont=none y=0.10-0.18
\textbf{光的衍射}．波长大衍射强，绕过障碍，$\lambda>$小孔尺\unclear{寸}时/障碍物尺寸时 明显衍射

光不沿直线传播了
%%%END
%%%BLOCK id=112-03 ch=15 sec=光的干涉·双缝 kind=knowledge alt=none cont=none y=0.14-0.42
\textbf{光的干涉}：同频(色)光、相位差不变、振动方向相同、只有同频率才能干涉

\struck{同\unclear{频}光衍射的叠加}\quad 同频光叠加

%FIG p112_a 0.213 0.168 0.39 0.267
\notefig[0.3]{figures/p112_a.png}

亮纹\quad $\Delta S=\dfrac{\lambda}{2}(2n)$\quad $n=0$ 第0亮纹 ...

暗纹\quad $\Delta S=\dfrac{\lambda}{2}(2n+1)$\quad \fbox{对一个\unclear{已知角}用余弦定理}

\[\Delta x=\frac{L'}{n-1}=\frac{L\lambda}{d}\] % CHECK: 分子疑为 L' 或 L 加一撇，原稿不清
$n$为$n$组条纹间距\quad 考察单位换算，数量级 % CHECK: “n为n组条纹间距”原文如此

双缝向下，$P$也向下（保同时性）

单色光\quad 红黑等距\quad 中亮边暗

复色光\quad 彩色等距\quad 中白边彩\quad 红光最外

绿光$+$蓝光$\Rightarrow$独立的蓝光条纹、绿光条纹
%%%END
%%%BLOCK id=112-04 ch=15 sec=光的干涉·薄膜 kind=knowledge alt=none cont=none y=0.42-0.55
\textbf{液体膜}(水/肥皂泡)

条纹间距\quad $x=\dfrac{\lambda}{2\tan\theta}=\dfrac{\lambda}{2\sin\theta}=\dfrac{\lambda}{2\theta}$\quad $\theta$很小\quad 液体膜顶角

%FIG p112_b 0.18 0.482 0.285 0.545
\notefig[0.25]{figures/p112_b.png}

前膜反射光和后膜反射光相干涉

光在人的一侧
%%%END
%%%BLOCK id=112-05 ch=15 sec=光的干涉·薄膜 kind=knowledge alt=none cont=none y=0.55-0.71
\textbf{气体膜}

% 图p112_c内含标注：窄 宽 窄（上方）；各处θ不同（下方）
%FIG p112_c 0.10 0.555 0.225 0.644
\notefig[0.3]{figures/p112_c.png}

% 图p112_d：圆环（牛顿环俯视），中心标“暗”
%FIG p112_d 0.17 0.643 0.27 0.71
\notefig[0.25]{figures/p112_d.png}

$x=\dfrac{\lambda}{2\tan\theta}$

$\theta$大条纹变密，$\theta$为气体膜顶角\quad 曲率

中间暗，往外为亮暗相间

% 图p112_e：楔形气体膜（带上下箭头）
%FIG p112_e 0.67 0.595 0.79 0.655
\notefig[0.25]{figures/p112_e.png}

% 图p112_f内含标注：人俯视条纹
%FIG p112_f 0.85 0.585 0.97 0.67
\notefig[0.25]{figures/p112_f.png}

\begin{tabular}{@{}cc@{\qquad}}
左凹 & 右凸 \\
同亮 & 同厚
\end{tabular}
%%%END
%%%BLOCK id=112-06 ch=15 sec=光的干涉·薄膜 kind=knowledge alt=none cont=none y=0.68-0.79
\textbf{固体膜}

\begin{itemize}
\item 增透膜(减反膜)\quad 增强透射光，减少反射光\\
  $2d=\dfrac{1}{2}\lambda\,(2n+1)$\qquad $d=\dfrac{1}{4}\lambda\cdot(2n+1)$
\item 增反膜(减透膜)\quad 增强反射光，减少透射光\\
  $2d=\dfrac{1}{2}\lambda\,(2n)$\qquad $d=\dfrac{1}{4}\lambda\cdot 2n$,
\end{itemize}

%FIG p112_g 0.555 0.72 0.645 0.785
\notefig[0.25]{figures/p112_g.png}
%%%END
%%%BLOCK id=112-07 ch=15 sec=光的衍射 kind=knowledge alt=none cont=none y=0.80-0.93
\textbf{单缝衍射、小孔衍射}，中间宽，两边暗．中宽边窄．中白边彩，红光最外

$\left\{\begin{tabular}{@{}l@{}}单缝越窄\\波长越长\end{tabular}\right.$ $\to$衍射越明显$\to$条纹间距大$\to$亮条纹范围大$\to$亮条纹亮度\unclear{低}

\textbf{圆盘衍射}(泊松亮斑)\quad 中窄边宽，中亮\quad 亮暗相间
%%%END
%%%BLOCK id=112-08 ch=15 sec=光的偏振 kind=knowledge alt=none cont=next y=0.95-1.00
% 本页底部被拍摄边缘截断，最后一行只露出上半部分
\textbf{偏振}：横波特有\quad 只通过平行光栅振动方向的光

偏振片可滤去反射光\quad 当偏振方向与反射光的偏振方向垂直时\unclear{…}(\unclear{横}波)\unclear{光…}

横波\quad 振动方向$\perp$传播方向

纵波\quad 振动方向$\parallel$传播方向
%%%END
%%%PAGEEND 112
